\section*{OBJECTS, MODELS AND STRUCTURES.}
\phantomsection
\addcontentsline{toc}{section}{OBJECTS, MODELS AND STRUCTURES.}

\begin{enumerate}
\def\labelenumi{\Alph{enumi})}
\tightlist
\item
  Objects and structures of mathematics. --- From Antiquity to the XIXth century, there is a common understanding as to what are the principal objects of the mathematicians; they are exactly those that are mentioned by Plato in the passage quoted earlier (p.~13); numbers, quantities and figures. If, at first, must be added the objects and phenomena which are the subject of Mechanics, Astronomy, Optics and Music, these ``mathematical'' disciplines are always clearly separated, amongst the Greeks, from Arithmetic and Geometry, and starting with the Renaissance they fairly quickly attain the rank of independent sciences.
\end{enumerate}

Whatever the philosophical nuances by which the conception of mathematical objects are coloured by such and such a mathematician or philosopher, there is at least one point on which there is unanimity: it is that these objects are given to us and it is not in our power to assign to them arbitrary properties, in the same way as a physicist cannot change a natural phenomenon. Truth to say, psychological reactions doubtless form one part of these views, reactions which are not for us to follow up, but which each mathematician knows when he wears himself out in vain efforts to seize upon a proof which seems to slip away endlessly. From there to assimilate this resistance to the obstacles that the world of senses sets in our path there is only a single step; and even today, more than one, who claims an intransigent formalism, would subscribe voluntarily, in his innermost being, to this admission of Hermite: ``I believe that numbers and the functions of Analysis are not the arbitrary result of our minds; I think that they exist outside of us, with the same character of necessity as the things of objective reality, and we meet them or discover them, and study them, as do the physicists, the chemists and the zoologists'' ({[}160{]}, v. II, p.~398).

There is no question in the classical conception of mathematics, of straying away from the study of numbers and figures; but this official doctrine, to which every mathematician believes himself bound to bring his verbal adherence, becomes no less bit by bit an intolerable burden, as new ideas accumulate. The embarrassment of the algebraists up against negative numbers does not cease until analytical Geometry gives it a convenient ``interpretation''; but, well into the XVIIIth century still, d' Alembert, discussing the question in the Encyclopédie ({[}75 a{]}, article NÉGATIF), loses his nerve suddenly after a column of fairly confused explanations, and is content to conclude that ``the rules of algebraic operations on negative quantities are generally assumed by everybody and perceived generally as exact, whatever idea is linked elsewhere to these quantities''. As for imaginary numbers, the scandal is far bigger; for if they are ``impossible'' roots and if (until about 1800) no way of ``interpreting'' them is seen, how can one, without contradiction, talk of these indefinable beings, and above all why introduce them? D'Alembert keeps here a prudent silence and does not even state these questions, no doubt because he realises that he would not be able to answer other than was done naively by A. Girard a century earlier ({[}129{]}, f.~22): ``It could be said: what use are these impossible solutions? I answer: for three things, for the certainty of the general rule, and that there is no other solution, and for its usefulness.''

In Analysis, the situation, in the XVIIth century, was no better. It is a happy circumstance that analytical Geometry was to appear, as if at a designated point, to give a ``representation'' in the shape of a geometric figure, of the great creation of the XVIIth century, the notion of a function, and so assist powerfully (with Fermat, Pascal or Barrow) at the birth of the infinitesimal Calculus (cf.~p.~193). But it is known on the other hand, to what philosophico-mathematical controversies the notions of the infinitely small and indivisible were going to give rise. And if d'Alembert is happier here, and recognises that in the ``metaphysics'' of the infinitesimal Calculus there is nothing other than the notion of limit ({[}75 a{]}, articles DIFFÉRENTIEL and LIMITE, and {[}75 b{]}), he is no more able than his contemporaries, to understand the real meaning of expansion in divergent series, and to explain the paradox of exact results obtained at the end of calculations with expressions deprived of any numerical interpretation. Finally, even in the domain of ``geometric certainty'', the Euclidean framework bursts: when Stirling, in 1717, does not shrink from saying that a certain curve has a ``double imaginary point at infinity'' ({[}299{]}, p.~93 of the new edition), he would certainly be in trouble trying to link such an ``object'' to commonly understood notions; and Poncelet who, at the beginning of the XIXth century, gives considerable development to such ideas in founding projective geometry (see p.~131), is content still to invoke as justification a ``principle of continuity'' which is altogether metaphysical.

One can conceive that, under these conditions (and at the exact moment when, paradoxically, the ``absolute truth'' of mathematics is proclaimed with the greatest force), the notion of proof seems to become more and more blurred during the XVIIIth century, since it is out of the question to set down, as did the Greeks, the notions on which reasoning is conducted, and their fundamental properties. The return towards rigour, which starts at the beginning of the XIXth century, brings some improvement to this state of things, but does not stop for all that the stream of new notions: one sees thus appearing in Algebra the imaginaries of Galois ({[}123{]}, p.~113-127), the ideal numbers of Kummer {[}188 b{]}, which are followed by vectors and quaternions, n-dimensional spaces, multivectors and tensors (see pp.~61 ff.), not to speak of Boolean algebra. No doubt great progress (which precisely allows the return to rigour, without losing any of the conquests of previous times) occurs with the possibility of giving ``models'' for these new notions in more classical terms: the ideal numbers or the imaginaries of Galois are interpreted by means of the theory of congruences (see pp.~82 ff.), n-dimensional geometry only appears (if so desired) as a pure language for stating results in algebra ``with n variables''; and for the classical imaginary numbers --- of which the geometric representation by the points of a plane (see pp.~161 ff.) marks the beginning of this flowering of Algebra --- there is soon the choice between the geometric ``model'' and an interpretation in terms of congruences (cf.~p.~82). But mathematicians begin at last to feel sharply that it is to fight against a natural slope along which their work is dragging them, and that it must be allowable, in mathematics, to reason about objects that have no sensible ``interpretation'': ``It is not of the essence of mathematics'', says Boole in 1854, ``to be occupied with the ideas of number and quantity'' ({[}29{]}, v. II, p.~13). \(^{30}\) The same preoccupation leads Grassmann, in his ``Ausdehnungslehre''

of 1844, to present his calculus in a form from which the notions of number or of geometric object are, at first, excluded. \(^{31}\) And, a bit later, Riemann, in his inaugural Lecture, takes care at the beginning not to speak of ``points'', but rather of ``determinations'' (Bestimmungsweise), in his description of ``multiplicities n times extended'', and underlines that, in such a multiplicity, the ``metric relations'' (Massverhältnisse) ``can only be studied for abstract quantities and can only be represented by formulae; under certain conditions, one can however decompose them into relations of which each taken in isolation is susceptible to a geometric representation, and in that way it is possible to express the results of the calculation in a geometric form'' ({[}259 a{]}, p.~276).

From that moment, the expansion of the axiomatic method is an accomplished fact. If, for a while longer, it is believed useful to control, when possible, the ``abstract'' results by geometric intuition, at least it is admitted that the ``classical'' objects are no longer the only ones that the mathematician can legitimately study. It is that --- precisely because of the multiple ``interpretations'' or ``models'' that are possible --- it has been recognised that the ``nature'' of mathematical objects is in the end secondary, and that it matters little, for example, that a result is presented as a theorem of ``pure'' geometry, or as a theorem of algebra by the means of analytical geometry. In other words, the essence of mathematics --- this elusive notion that could until then only have been expressed by vague names such as ``general rule'' or ``metaphysics'' --- appears as the study of relationships between objects that are only (voluntarily) known and described by some of their properties, precisely those that are put as axioms at the foundations of their theory. It is this that had already been clearly seen by Boole in 1847, when he wrote that mathematics deals with ``operations considered in themselves, independently of the diverse objects to which they can be applied'' ({[}29{]}, v. I, p.~3). Hankel, in 1867, inaugurating the axiomatisation of algebra, defends a mathematics that is ``purely intellectual, a pure theory of forms, which has as its purpose, not the combination of quantities, or of their images, the numbers, but objects of thought (``Gedankendinge'') to which may correspond effective objects or relations, even though such a correspondence is not necessary'' ({[}146{]}, p.10). Cantor, in 1883, echoes this claim of a ``free mathematics'' by proclaiming that ``mathematics is entirely free in its development, and its concepts are only linked by the necessity of being consistent, and are co-ordinated with concepts introduced previously by means of precise definitions'' ({[}47{]}, p.182). Finally, the revision of Euclidean geometry succeeds in spreading and popularising these ideas. Pasch himself, although still attached to a certain ``reality'' of geometric objects, recognises that geometry is in fact independent of their significance, and consists purely in the study of their relations ({[}245{]}, p.90); a concept that Hilbert pushes to its logical conclusion in underlining that even the names of the basic notions of a mathematical theory can be chosen arbitrarily, \(^{32}\) and that Poincaré expresses by saying that the axioms are ``disguised definitions'', thus completely reversing the scholastic point of view.

It would thus be tempting to say that the modern notion of ``structure'' is attained in substance around 1900; in fact it will need still another thirty years of apprenticeship before it appears in all its glory. It is no doubt not difficult to recognise structures of the same kind when they have a fairly simple nature; for the group structure, for example, this stage was reached already by the middle of the XIXth century. But at the same time, Hankel can be seen fighting --- without quite succeeding --- to bring out the general ideas of field and extension, that he does not succeed in expressing except in the form of a ``principle of permanence'' which is semi metaphysical {[}146{]}, and which will be only formulated in a definitive form by Steinitz {[}294 a{]} 40 years later. Above all it has been fairly difficult, in this matter, to free oneself from the impression that mathematical objects are ``given'' to us with their structure; only a fairly long use of functional Analysis has been able to familiarise modern mathematicians with the idea that, for example, there are several ``natural'' topologies on the rational numbers, and several measures on the number line. With this disassociation the transition to general definition of structures was finally realised.

\begin{enumerate}
\def\labelenumi{\Alph{enumi})}
\setcounter{enumi}{1}
\item
  Models and isomorphisms. --- The intervention of the notion of ``model'' or ``interpretation'' of one mathematical theory by means of another will have been noticed on several occasions. That is not a recent idea, and no doubt what can be seen here is a manifestation ceaselessly recurring of a deep feeling for the unity of the various ``mathematical sciences''. If one can take as authentic the traditional maxim ``All things are numbers'' of the first Pythagoreans, that can be considered as the trace of a first attempt to bring back the geometry and algebra of the times to arithmetic. Although the discovery of the irrationals seemed to close for ever that route, the reaction to which it gave birth in Greek mathematics was a second attempt at synthesis this time taking geometry as the basis, and absorbing among others the solution methods for algebraic equations inherited from the Babylonians. \(^{33}\) It is known that this conception was to remain until the fundamental reform of R. Bombelli and of Descartes, assimilating every measure of quantity to a measure of length (in other words, to a real number; cf.~p.~151). But with the creation of analytical geometry by Descartes and Fermat, the tendency is again reversed, and a much tighter fusion of geometry and algebra is obtained, but this time to the benefit of algebra. Further, elsewhere, Descartes goes further and conceives of the essential unity of ``all these sciences that are commonly called Mathematics \ldots{} Even though their subjects are different'', he says ``they still do not fail to converge, in that they consider nothing other than the various relations or ratios that are to be found there'' ({[}85 a{]}, v. VI, pp.~19-20). \(^{34}\) However, this point of view only tended to make Algebra the fundamental mathematical science; a conclusion against which Leibniz protests vigorously, who also himself, as has been seen, conceived of a ``universal Mathematics'', but on a much vaster scale and already quite close to modern ideas. Making precise the ``accord'' of which Descartes spoke, he glimpses, in fact, for the first time, the general notion of isomorphism (which he calls ``similitude''), and the possibility of ``identifying'' relations or operations that are isomorphic; he gives examples addition and multiplication ({[}69 a{]}, pp.~301-303). But these audacious views remained without echo amongst his contemporaries, and one must await the expansion of Algebra which takes place around the middle of the XIXth century (see pp.~51 ff.) to see the beginnings of the realisation of the Leibnizian dreams. We have already underlined that it is at that moment that the ``models'' multiply and that it becomes usual to go from one theory to another by a simple change of language; the most striking example of this is perhaps duality in projective geometry (see p.~132), where the practice, frequent at that time, of printing face to face, in two columns, the theorems ``dual'' one to the other, doubtless plays a great part in the full realisation of the notion of isomorphism. From a more technical point of view, certainly the notion of isomorphic groups for abelian groups is known to Gauss, for groups of permutations to Galois (cf.~pp.~51 ff.); it is acquired in a general way for any group whatsoever around the middle of the XIXth century. \(^{35}\) Following this, with each new axiomatic theory, it is a natural development to define a notion of isomorphism; but it is only with the modern notion of structure that it was finally recognised that every structure carries within itself a notion of isomorphism, and that it is not necessary to give a special definition of it for each type of structure.
\item
  The arithmetisation of classical mathematics. --- The use, more and more widespread, of the notion of ``model'' was also going to allow the realisation in the XIXth century of the unification of mathematics dreamed of by the Pythagoreans. At the beginning of the century, whole numbers and continuous quantities seemed as irreconcilable as ever in antiquity; real numbers remain linked to the notion of geometric quantity (at least to that of length), and it is to this latter that appeal had been made for the ``models'' of negative numbers and imaginary numbers. Even rational numbers were traditionally attached to the idea of the splitting of a quantity into equal parts; only the whole numbers remained apart, as ``exclusively the product of our minds'' as is said by Gauss in 1832, putting them in opposition to the notion of space ({[}124 a{]}, v. VIII, p.~201). The first efforts to bring together Arithmetic and Analysis were made at first with rational numbers (positive and negative) and are due to Martin Ohm (1822); they were taken up again around 1860 by several authors, notably Grassmann, Hankel and Weierstrass (in his unpublished lectures); it is to this latter that it appears is due the idea of obtaining a ``model'' of the positive rational numbers or of whole negative numbers by considering classes of pairs of whole numbers. But the most important step remained to be taken, namely to find a ``model'' of irrational numbers in the theory of rational numbers; around 1870, it had become an urgent problem, in view of the necessity, after the discovery of ``pathological'' phenomena in Analysis, to eliminate every trace of geometric intuition in the vague notion of ``quantity'' in the definition of real numbers. It is known that this problem was solved around this time, almost simultaneously by Cantor, Dedekind, Méray and Weierstrass, and using fairly different methods (see p.~155).
\end{enumerate}

From this moment, whole numbers became the foundation of all classical mathematics. Further, ``models'' based on Arithmetic acquired still more importance with the extension of the axiomatic method and the conception of mathematical objects as free creations of the mind. There remained in fact one restriction to this freedom claimed by Cantor, the question of ``existence'' which had already preoccupied the Greeks, and which arose here in a much more immediate way, precisely since all appeals to an intuitive representation had now been abandoned. We will see later (pp.~36 ff.) of what a philosophico-mathematical maelstrom the notion of ``existence'' was going to be the centre in the first years of the XXth century. But in the XIXth century that had not been reached, and to prove the existence of a mathematical object having given properties, it was simply, as for Euclid, ``constructing'' an object with given properties. That was precisely what arithmetic ``models'' were for: once the real numbers were ``interpreted'' in terms of whole numbers, complex numbers and Euclidean geometry were also, thanks to analytical Geometry, and it was the same for all new algebraic objects introduced since the beginning of the century; finally --- a discovery that had a great effect --- Beltrami and Klein had even obtained Euclidean ``models'' of the non-Euclidean geometries of Lobatschevsky and of Riemann (see p.~134), and in consequence ``arithmetised'' (and thereby completely justified) these theories which had at first sight aroused such distrust.

\begin{enumerate}
\def\labelenumi{\Alph{enumi})}
\setcounter{enumi}{3}
\tightlist
\item
  The axiomatisation of arithmetic.---It was in line with this evolution that a subsequent turn towards the foundations of arithmetic itself was made, and indeed that is what can be seen around 1880. It appears that before the XIXth century no attempt had been made to define addition and multiplication of whole numbers except by a direct appeal to intuition; Leibniz is the only one who, faithful to his principles, warns expressly that ``truths'' as ``obvious'' as \(2 + 2 = 4\) are no less susceptible to proof if one thinks about the definitions of the numbers involved ({[}198 b{]}, v. IV, p.~403; cf.~{[}69 a{]}, p.~203); and he did not consider at all that commutativity of addition and multiplication were intrinsic. \(^{36}\) But he does not take his thoughts on this subject any further, and about the middle of the XIXth century, still no progress in this direction had been made: Weierstrass himself, whose lectures contributed considerably to extend the ``arithmeticising'' point of view, does not seem to have felt the need for a logical clarification of the theory of whole numbers. The first steps in this direction seem to be due to Grassmann, who, in 1861 ({[}134{]}, v. II \(_{2}\) , p.~295) gives a definition of addition and multiplication for whole numbers, and proves their fundamental properties (commutativity, associativity, distributativity) using only the operation \(x \mapsto x + 1\) and the principle of induction. This latter had been clearly conceived and used for the first time in the XVIIth century by B. Pascal ({[}244{]}, v. III, p.~456) \(^{37}\) --- even though one can find in Antiquity some more or less conscious applications of it --- and was currently used by mathematicians since the second half of the XVIIth century. But it is only in 1888 that Dedekind ({[}79{]}, v. III, pp.359-361) stated a complete system of axioms for arithmetic (a system repeated 3 years later by Peano and usually known by his name {[}246 c{]}), which contained in particular a precise formulation of the principle of induction (that Grassmann uses yet without stating it explicitly).
\end{enumerate}

With this axiomatisation, it seemed that the definitive foundations of mathematics had been reached. In fact, at the very moment when the axioms of arithmetic were being clearly formulated, these same, for many mathematicians (starting with Dedekind and Peano themselves) were already deprived of this role of primordial science, in favour of the latest arrival among the theories of mathematics, the theory of sets; and the controversies that were going to unravel around the notion of whole number cannot be isolated from the great ``crisis of the foundations'' of the years 1900-1930.
% semantic-fragments: legacy-input-seam
\relax{}
